% !TEX spellcheck = en_US

\documentclass[11pt]{amsart}



% ~~~~~~~~~~~~~~~~~~~~~~~~~~~~~~~~~~~~~~~ V lib
% Note that hblibCommon is loaded by hblibArticle
\usepackage{../../cis000_latex/lib/hblibArticle}
\usepackage{../../cis000_latex/lib/hblibDatastructures}
\usepackage{../../cis000_latex/lib/hblibCodeoutput}
% ~~~~~~~~~~~~~~~~~~~~~~~~~~~~~~~~~~~~~~~ A lib



% ~~~~~~~~~~~~~~~~~~~~~~~~~~~~~~~~~~~~~~~ V specific
% ~~~~~~~~~~~~~~~~~~~~~~~~~~~~~~~~~~~~~~~ A specific



% ~~~~~~~~~~~~~~~~~~~~~~~~~~~~~~~~~~~~~~~ A titleParam
\newcommand{\cisno}{111}	% 111		
\newcommand{\courseName}{Fundamentals of Software Development}	% 111		
\newcommand{\semesterID}{2026-3}	% 2026-3	
\newcommand{\semesterName}{Fall}	% Fall	
\newcommand{\courseInstructors}{Bingol$^{*}$, Alpaydin, Demirer, Erkovan}		

\newcommand{\examDate}{2026-11-22}	% 2026-11-22	
\newcommand{\examNo}{1m1}	% 1m1	
\newcommand{\examName}{Midterm-1}	% 1m1	
\newcommand{\questionNo}{Q2}	% Q1	
\newcommand{\groupNo}{g1}	% g1	

\newcommand{\topic}{Doubly Linked List}	% Array Library
\newcommand{\topicShort}{DLL}		% ArrayLib
\newcommand{\questionAuthor}{Haluk Bingol}		% Haluk Bingol
\newcommand{\questionAuthorShort}{Bingol}		% Bingol
% ~~~~~~~~~~~~~~~~~~~~~~~~~~~~~~~~~~~~~~~ A titleParam



% ~~~~~~~~~~~~~~~~~~~~~~~~~~~~~~~~~~~~~~~ V frontExam
% defines: \courseID, \questionID, \hStudentInfo, \refStudentInfo,
% \textUpload
% Common title format for exam
% Common warning
% !TEX spellcheck = en_US

% ~~~~~~~~~~~~~~~~~~~~~~~~~~~~~~~~~~~~~~~ V title



% ~~~~~~~~~~~~~~~~~~~~~~~~~~~~~~~~~~~~~~~ V derived
% do not change
\newcommand{\courseID}{\textsc{cis}\cisno}
\newcommand{\questionID}{\courseID-\semesterID-\examNo-\questionNo-\groupNo}
% ~~~~~~~~~~~~~~~~~~~~~~~~~~~~~~~~~~~~~~~ A derived


% ~~~~~~~~~~~~~~~~~~~~~~~~~~~~~~~~~~~~~~~ V cmdApp
% Commands defined in appendix
% \hStudentInfo;  \refStudentInfo
% \hBelowAboveMarker;   \refBelowAboveMarker
 % Commands defined in appendix


\newcommand{\hStudentInfo}{ \hCode{StudentInfo}}
\newcommand{\hBelowAboveMarker}{ \hCode{BelowAboveMarker}}

\newcommand{\refStudentInfo}{\ref{sec:app:studentInfo}}
\newcommand{\refBelowAbove}{\ref{sec:app:belowAboveMarker}}

\newcommand{\textUpload}{%
Upload your work to \YULearn. 
See \refsec{sec:app:upload}.
}

% ~~~~~~~~~~~~~~~~~~~~~~~~~~~~~~~~~~~~~~~ A cmdApp


% ~~~~~~~~~~~~~~~~~~~~~~~~~~~~~~~~~~~~~~~ V title
% Common title format for exam
% !TEX spellcheck = en_US

% ~~~~~~~~~~~~~~~~~~~~~~~~~~~~~~~~~~~~~~~ V title
% Common title format for exam question

\title[\courseID-\examNo: \topicShort]{
% ~~~~~~~~~~~~~~~~~~~~~~~~~~~~~~~~~~~~~~~ V box
	\rule[1ex]{\textwidth}{.1pt}
	\begin{tabular*}{\textwidth}{l @{\extracolsep{\fill}} r}
		\courseID\ \courseName
			&\semesterName\ \semesterID\\
		\courseInstructors
			&\examDate\\
		\examName
			&\questionID\\
	\end{tabular*}
	\rule[1ex]{\textwidth}{.1pt}
	~\\
	\vspace{1cm}
% ~~~~~~~~~~~~~~~~~~~~~~~~~~~~~~~~~~~~~~~ A box
	{\Large 
		\questionNo\ . \topic
	}
}

\author{
	\questionAuthor\\
	\hbTimeStamp
}

\begin{document}

\maketitle

\tableofcontents
% ~~~~~~~~~~~~~~~~~~~~~~~~~~~~~~~~~~~~~~~ A title

% ~~~~~~~~~~~~~~~~~~~~~~~~~~~~~~~~~~~~~~~ A title


% ~~~~~~~~~~~~~~~~~~~~~~~~~~~~~~~~~~~~~~~ V title
% Common warning
% !TEX spellcheck = en_US

% ~~~~~~~~~~~~~~~~~~~~~~~~~~~~~~~~~~~~~~~ V warning
% Common title format for exam question
\begin{hBoxWarning}
Please
\begin{itemize}
	
	\item 
	read \refapp{sec:appExaminfo} Exam Info.
	
	\item
	populate \hStudentInfo.
	See \ref{sec:app:studentInfo}.
	
	\item
	modify only within ``below'' and ``above'' markers.
	See \ref{sec:app:belowAboveMarker}.

	\item
	read warning about AutoGrader.
	See \ref{sec:app:warnAutoGrader}.
	
\end{itemize}
\end{hBoxWarning}
% ~~~~~~~~~~~~~~~~~~~~~~~~~~~~~~~~~~~~~~~ A warning

% ~~~~~~~~~~~~~~~~~~~~~~~~~~~~~~~~~~~~~~~ A title


% ~~~~~~~~~~~~~~~~~~~~~~~~~~~~~~~~~~~~~~~ A title

% ~~~~~~~~~~~~~~~~~~~~~~~~~~~~~~~~~~~~~~~ A frontExam



% ~~~~~~~~~~~~~~~~~~~~~~~~~~~~~~~~~~~~~~~ V
\section{Introduction}
\label{sec:intro}

Some intro here such as the following.

You may use \YULearn, 
student info \refStudentInfo\ 
or ``below-above'' markers \refBelowAbove.

A list is given by a reference to its \textbf{first node}, called \hCode{list}.
We write \hCode{list -> [1, 5, 4]} for the list in \reffig{fig:list}.
The \textbf{empty list} is \hCode{null}.

% +++++++++++++++++++++++++++++++++++++++ V
%: -Figure. [1, 5, 4]
\begin{figure}[h]
\centering
\begin{tikzpicture}
	\hDLList{list}{0}{0}{1,5,4}{0}
\end{tikzpicture}
\caption{\hCode{list -> [1, 5, 4]}. 
	Each node is composed of  \hCode{prev}, \hCode{data}, and \hCode{next} fields.
	\textcolor{blue!70!black}{blue links} represent \hCode{next},
	\textcolor{red!70!black}{red links} represent \hCode{prev} links.}
	\label{fig:list}
\end{figure}
% +++++++++++++++++++++++++++++++++++++++ A


Use libraries in \hCode{hblibDatastructures}.
% ~~~~~~~~~~~~~~~~~~~~~~~~~~~~~~~~~~~~~~~ V
Visual representation of a stack after
\begin{javacode}
StringStack s = new StringStack(6);

s.push("a");
s.push("b");
s.push("c");
\end{javacode}
is the following:
% +++++++++++++++++++++++++++++++++++++++ V
%: -Figure.
\begin{center}
\begin{tikzpicture} 
	\hStack{0}{6}{2}{a,b,c}[0.7][0.6] 
\end{tikzpicture}
\end{center}
% +++++++++++++++++++++++++++++++++++++++ A

% ~~~~~~~~~~~~~~~~~~~~~~~~~~~~~~~~~~~~~~~ A



% ~~~~~~~~~~~~~~~~~~~~~~~~~~~~~~~~~~~~~~~ V
% switch from numbers to Q
\hSectionQuestion
% ~~~~~~~~~~~~~~~~~~~~~~~~~~~~~~~~~~~~~~~ A



% +++++++++++++++++++++++++++++++++++++++ V modifyMe
%: -Table.
\begin{table}[h]
\centering\small
\begin{tabular}{l}
	\toprule
	ModifyMe\\
	\midrule
	\hCode{DoublyLinkedList.java}\\
	\bottomrule
\end{tabular}
\caption{
	List of files that you are \textbf{allowed to modify}. 
	AutoGrader will use the files in this list only.
	\textbf{Do not modify any other file.}
	\label{tbl:modifyMe}
}
\end{table}
% +++++++++++++++++++++++++++++++++++++++ A modifyMe



% ~~~~~~~~~~~~~~~~~~~~~~~~~~~~~~~~~~~~~~~ V Q1
\subsection{\hCode{sumOfNeighbors} method}
\label{sec:sumOfNeighbors}

There is no rule in question body.
This is one possible way.
Feel free to change it.



% ~~~~~~~~~~~~~~~~~~~~~~~~~~~~~~~~~~~~~~~
\subsubsection{Signature}

% +++++++++++++++++++++++++++++++++++++++ V
%: -Code.
\begin{javacode}
public static int[] sumOfNeighbors(int[] arr)
\end{javacode}
% +++++++++++++++++++++++++++++++++++++++ A



% ~~~~~~~~~~~~~~~~~~~~~~~~~~~~~~~~~~~~~~~
\subsubsection{Function} 

Description of the function.



% ~~~~~~~~~~~~~~~~~~~~~~~~~~~~~~~~~~~~~~~
\subsubsection{Rules}

Some rules.



% ~~~~~~~~~~~~~~~~~~~~~~~~~~~~~~~~~~~~~~~
\subsubsection{Example}

Some examples.



% ~~~~~~~~~~~~~~~~~~~~~~~~~~~~~~~~~~~~~~~
\subsubsection{Constraints}

Some constraints.

\textbf{Hint.} 
Some hint.


% ~~~~~~~~~~~~~~~~~~~~~~~~~~~~~~~~~~~~~~~ A Q1




% ~~~~~~~~~~~~~~~~~~~~~~~~~~~~~~~~~~~~~~~ V
% switch from Q back to numbers
\hNumbersArabic 
% ~~~~~~~~~~~~~~~~~~~~~~~~~~~~~~~~~~~~~~~ A


% ~~~~~~~~~~~~~~~~~~~~~~~~~~~~~~~~~~~~~~~ V testing
\section{Testing}

Use \hCode{TS\_xxx\_jUnit} to test your code.
You can also use \hCode{xxx\_Test} if there is any.
% ~~~~~~~~~~~~~~~~~~~~~~~~~~~~~~~~~~~~~~~ A testing




% ~~~~~~~~~~~~~~~~~~~~~~~~~~~~~~~~~~~~~~~ V
\section{ToDo}\label{sec:todo}

\begin{enumerate}
	\item Fill in your personal information at the beginning of \hCode{ArrayLib.java}.
	\item Complete the \hCode{sumOfNeighbors} method.
	\item Complete the \hCode{productOfOdds} method.
	\item Complete the \hCode{selectAtMost} method.
	\item Write your code only between the \hCode{below} and \hCode{above} marks.
	\item Run \hCode{TS\_ArrayLib\_jUnit} to verify your implementation.
	
	\item \textUpload 
\end{enumerate}
% ~~~~~~~~~~~~~~~~~~~~~~~~~~~~~~~~~~~~~~~ A



% ~~~~~~~~~~~~~~~~~~~~~~~~~~~~~~~~~~~~~~~ V \appendix
% insert appendix
% Exam info

\newpage
\appendix


% ~~~~~~~~~~~~~~~~~~~~~~~~~~~~~~~~~~~~~~~ V
\section{Exam}
\input{../cis000_latex/input/in-ExamInfo-Intro}
\input{../cis000_latex/input/in-Period1-BeforeExamStarts.tex}
\input{../cis000_latex/input/in-Period2-DuringExam.tex}
\input{../cis000_latex/input/in-Period3-AtTheEndOfTheExam.tex}
\input{../cis000_latex/input/in-WarningAboutTheAutoGrader}
\input{../cis000_latex/input/in-ExamTiming}	
\input{../cis000_latex/input/in-Objection}
% ~~~~~~~~~~~~~~~~~~~~~~~~~~~~~~~~~~~~~~~ A

% ~~~~~~~~~~~~~~~~~~~~~~~~~~~~~~~~~~~~~~~ A \appendix


% ~~~~~~~~~~~~~~~~~~~~~~~~~~~~~~~~~~~~~~~ 
\end{document}
